\exercisesection{}

Unless stated otherwise, in the exercises of this section, E denotes a complex Hilbert space.

\begin{enumerate}
\def\labelenumi{\arabic{enumi})}
\tightlist
\item
  Let \(u\) be an endomorphism of E whose image is closed in E. Let \(q\) be an orthogonal projector of E and let \(v\) be the endomorphism \(x \mapsto q(u(x))\) of the image of E.
\end{enumerate}

\begin{enumerate}
\def\labelenumi{\alph{enumi})}
\item
  The adjoint of \(v\) is the endomorphism \(x \mapsto q(u^{*}(x))\) of the image of E.
\item
  If \(u\) is Hermitian, then \(v\) is Hermitian.
\item
  It is possible that \(u\) is normal and that \(v\) is not.
\end{enumerate}

\begin{enumerate}
\def\labelenumi{\arabic{enumi})}
\setcounter{enumi}{1}
\item
  Prove Theorem 1 of IV, p. 149 by using directly the spectral theorem (cor. of th. 1 of IV, p. 193).
\item
  Let \(S \subset \mathbf{C}\) be a compact subset of \(\mathbf{C}\). Let \(\varphi \colon \mathcal{C}(S) \to \mathcal{L}(E)\) be a continuous unital morphism of involutive algebras. There exists a unique normal endomorphism \(u\) of \(E\) whose spectrum is \(S\) and which satisfies \(\varphi(f) = f(u)\) for all \(f \in \mathcal{L}^{\infty}(S)\) .
\item
  Let \(u \in \mathcal{L}(\mathrm{E})\) be normal. The convex hull of the spectrum of \(u\) is the closure in \(\mathbf{C}\) of the numerical range \(\iota(u)\) of \(u\) (cf. I, p. 135, def. 2).
\item
  Let \(u \in \mathcal{L}(\mathrm{E})\) be such that \(\| u \| \leqslant 1\) . For every \(\lambda \in \mathbf{C}\) , let \(p_{\lambda}\) be the orthogonal projector of E onto the eigenspace of \(u\) corresponding to \(\lambda\) . For all \(x\) and \(y\) in E, we have
\end{enumerate}

\[
\lim_{N\to +\infty}\frac{1}{N}\sum_{n = 1}^{N}|\langle x \mid u^{n}(y)\rangle |^{2} = \sum_{\substack{\lambda \in \mathbf{C}\\ |\lambda | = 1}}|\langle p_{\lambda}(y) \mid p_{\lambda}(x)\rangle^{2}|\\ = \sum_{\substack{\lambda \in \mathbf{C}\\ |\lambda | = 1}}|\langle y \mid p_{\lambda}(x)\rangle |^{2}
\]

(one may use Exercise 64 of II, p. 300).

\begin{enumerate}
\def\labelenumi{\arabic{enumi})}
\setcounter{enumi}{5}
\tightlist
\item
  Let E be a pre-Hilbert space, let F be a Banach space, and let \(u: E \rightarrow F\) be a linear map.
\end{enumerate}

\begin{enumerate}
\def\labelenumi{\alph{enumi})}
\item
  Prove that u is a compact linear map if and only if \(\|u(e_{n})\|\to0\) for every orthonormal sequence \((e_{n})\) of E. (To prove that this condition is sufficient, use TG, IX, p. 20, prop. 14.)
\item
  Assume that E is a separable Hilbert space and that F = E. For there to exist an orthonormal basis \((e_{n})\) of E such that \(\|u(e_{n})\| \to 0\) it is necessary and sufficient that u not be a Fredholm operator.
\item
  We assume that E is a Hilbert space and that we have \(\langle u(e_{n}), e_{n} \rangle \to 0\) for every orthonormal sequence \((e_{n})\) of H. Prove that u is a compact linear map.
\end{enumerate}

\begin{enumerate}
\def\labelenumi{\arabic{enumi})}
\setcounter{enumi}{6}
\tightlist
\item
  Let \(u\) and \(v\) be continuous endomorphisms of E.
\end{enumerate}

\begin{enumerate}
\def\labelenumi{\alph{enumi})}
\item
  Prove that \(\operatorname{Im}(u) + \operatorname{Im}(v) = \operatorname{Im}((u \circ u^* + v \circ v^*)^{1/2})\) .
\item
  If \(u\) is positive and has closed image; prove that \(\operatorname{Im}(u) = \operatorname{Im}(u^{1/2})\) .
\item
  If \(u, v\) are positive with closed image, and if \(u + v\) is also closed-image, we have \(\operatorname{Im}(u + v) = \operatorname{Im}(u) + \operatorname{Im}(v)\) .
\end{enumerate}

\begin{enumerate}
\def\labelenumi{\arabic{enumi})}
\setcounter{enumi}{7}
\tightlist
\item
  Let F be a complex Hilbert space and let \(u \colon \mathrm{E} \to \mathrm{F}\) be a continuous linear map with closed image.
\end{enumerate}

\begin{enumerate}
\def\labelenumi{\alph{enumi})}
\tightlist
\item
  Prove that there exists a unique continuous linear map \(v \in \mathcal{L}(\mathrm{F};\mathrm{E})\) satisfying the following conditions:
\end{enumerate}

\begin{enumerate}
\def\labelenumi{(\roman{enumi})}
\item
  \(u \circ v \circ u = u\) and \(v \circ u \circ v = v\) ;
\item
  \(u \circ v\) and \(v \circ u\) are self-adjoint.
\end{enumerate}

We denote this map by \(u^{\dagger}\) (Moore--Penrose pseudoinverse).

\begin{enumerate}
\def\labelenumi{\alph{enumi})}
\setcounter{enumi}{1}
\tightlist
\item
  Prove that \(\operatorname{Im}(u) = \operatorname{Im}(u \circ u^{\dagger})\). Show that for \(x \in E\) and \(y \in \operatorname{Im}(u)\) , the conditions \(u(x) = y\) and \(x - u^{\dagger}(y) \in \operatorname{Im}(1_{\operatorname{E}} - u^{\dagger} \circ u)\) are equivalent.
\end{enumerate}

\begin{enumerate}
\def\labelenumi{\arabic{enumi})}
\setcounter{enumi}{8}
\tightlist
\item
  Let \(u\) and \(v\) be endomorphisms of E with closed images. If \(u + v\) also has closed image, we set \(u: v = u \circ (u + v)^{\dagger} \circ v\) , where \((u + v)^{\dagger}\) is the Moore--Penrose pseudoinverse of \(u + v\) .
\end{enumerate}

\begin{enumerate}
\def\labelenumi{\alph{enumi})}
\item
  Prove that \(u:v\) is self-adjoint.
\item
  One has \(u:v=v:u\) and \(\operatorname{Im}(u:v)=\operatorname{Im}(u)\cap\operatorname{Im}(v)\) .
\item
  Assume that u and v are orthogonal projectors. Prove that 2u : v is the orthogonal projector with image \(\operatorname{Im}(u) \cap \operatorname{Im}(v)\) .
\end{enumerate}

\begin{enumerate}
\def\labelenumi{\arabic{enumi})}
\setcounter{enumi}{9}
\tightlist
\item
  A continuous endomorphism \(u\) of E is said to be hyponormal if one has \(u^{*} \circ u \geqslant u \circ u^{*}\), and is said to be paranormal if \(\|u^{2}(x)\|\|x\| \geqslant \|u(x)\|^{2}\) for all \(x \in E\) .
\end{enumerate}

\begin{enumerate}
\def\labelenumi{\alph{enumi})}
\item
  If \(u\) is paranormal, then \(u^n\) is paranormal for every \(n \in \mathbf{N}\) .
\item
  If u is paranormal and invertible, then \(u^{-1}\) is paranormal.
\item
  A hyponormal endomorphism is paranormal.
\item
  Let u be a paranormal endomorphism whose spectrum is contained in the unit circle; then u is unitary.
\item
  If \(u\) is a paranormal endomorphism, its norm \(\| u\|\) is equal to its spectral radius \(\varrho(u)\) .
\item
  We assume that u is paranormal.
\end{enumerate}

\begin{enumerate}
\def\labelenumi{\arabic{enumi})}
\setcounter{enumi}{10}
\tightlist
\item
  Let u be a continuous endomorphism of E.
\end{enumerate}

\begin{enumerate}
\def\labelenumi{\alph{enumi})}
\tightlist
\item
  For \(u\) to be paranormal, it is necessary and sufficient that
\end{enumerate}

\[
(u ^ {*}) ^ {2} \circ u ^ {2} - 2 t (u ^ {*} \circ u) + t ^ {2} 1 _ {\mathrm{E}} \geqslant 0
\]

for every real number t \textgreater{} 0.

\begin{enumerate}
\def\labelenumi{\alph{enumi})}
\setcounter{enumi}{1}
\tightlist
\item
  For \(u\) to be paranormal, it is necessary and sufficient that
\end{enumerate}

\[
p \circ q ^ {2} \circ p - 2 t p ^ {2} + t 1 _ {\mathrm{E}} \geqslant 0
\]

for all \(t > 0\) , where \(p = (u \circ u^{*})^{1/2}\) and \(q = (u^{*} \circ u)^{1/2}\) .

\begin{enumerate}
\def\labelenumi{\arabic{enumi})}
\setcounter{enumi}{11}
\tightlist
\item
  \begin{enumerate}
  \def\labelenumii{\alph{enumii})}
  \tightlist
  \item
    Let u and v be continuous positive endomorphisms of E. Suppose that for every t \textgreater{} 0, we have
  \end{enumerate}
\end{enumerate}

\[
2 t u ^ {2} \circ (u ^ {2} + t ^ {2} 1 _ {\mathrm{E}}) ^ {- 1} \leqslant v \leqslant \frac {1}{2 t} (u ^ {2} + t ^ {2} 1 _ {\mathrm{E}}).
\]

Prove that u = v.

\begin{enumerate}
\def\labelenumi{\alph{enumi})}
\setcounter{enumi}{1}
\tightlist
\item
  For \(u\) to be a normal endomorphism, it is necessary and sufficient that \(u\) and \(u^*\) be paranormal and have the same kernel.
\end{enumerate}

\begin{enumerate}
\def\labelenumi{\arabic{enumi})}
\setcounter{enumi}{12}
\tightlist
\item
  \begin{enumerate}
  \def\labelenumii{\alph{enumii})}
  \tightlist
  \item
    Let A be an infinite set of unit vectors of E which converges weakly to 0 along the filter of complements of finite sets. Prove that there exists a sequence \((x_{n})_{n\geqslant 1}\) with values in A and an orthonormal sequence \((e_n)_{n\geqslant 1}\) in E such that \(\| x_n - e_n\| \leqslant 1 / n\) for all \(n\geqslant 1\) .
  \end{enumerate}
\end{enumerate}

\begin{enumerate}
\def\labelenumi{\alph{enumi})}
\setcounter{enumi}{1}
\item
  Let u be an endomorphism of E, not necessarily continuous. Suppose that for every orthonormal sequence \((e_{n})\) in E, the sequence \((u(e_{n}))\) converges weakly to 0. Prove that u is continuous and compact.
\item
  Let u be a continuous endomorphism of E. For there to exist an orthonormal basis \((e_{n})\) of E such that the sequence \((u(e_{n}))\) converges weakly to 0, it is necessary and sufficient that for every endomorphism v of E, \(v \circ u - 1_{E}\) is not compact.
\end{enumerate}

\begin{enumerate}
\def\labelenumi{\arabic{enumi})}
\setcounter{enumi}{13}
\tightlist
\item
  Let X be a locally compact space, \(\mu\) a positive bounded measure on X such that \(\mu(\mathrm{X}) = 1\) , and let \(f: X \to X\) be a \(\mu\)-measurable map such that \(f(\mu) = \mu\) . We denote by \(u_{f}\) the unitary endomorphism of \(\mathrm{L}_{\mathbf{C}}^{2}(\mathrm{X}, \mu)\) defined by \(\varphi \mapsto f \circ \varphi\) (cf. exercise 16 of I, p. 190). For every \(n \in N\), we denote by \(f^{n} = f^{\circ n}\) (the \(n\)-fold composition), and for every subset A of X, we denote \(f^{-n}(\mathrm{A}) = g^{-1}(\mathrm{A})\) with \(g = f^{n}\) .
\end{enumerate}

We say that \(f\) is weakly mixing relative to \(\mu\) if, for all \(\mu\)-measurable subsets A and B of X, we have

\[
\lim _ {N \rightarrow + \infty} \frac {1}{N} \sum_ {n = 0} ^ {N} | \mu (A \cap f ^ {- n} (B) - \mu (A) \mu (B) | = 0.
\]

\begin{enumerate}
\def\labelenumi{\alph{enumi})}
\tightlist
\item
  The following conditions are equivalent:
\end{enumerate}

\begin{enumerate}
\def\labelenumi{(\roman{enumi})}
\item
  The map f is weakly mixing relative to \(\mu\) ;
\item
  For all \(\varphi_{1},\varphi_{2} \in \mathrm{L}_{\mathbf{C}}^{2}(\mathrm{X}, \mu)\) , we have
\end{enumerate}

\[
\lim _ {N \to + \infty} \frac {1}{N} \sum_ {n = 0} ^ {N} | \langle \varphi_ {1} | u _ {f} ^ {n} (\varphi_ {2}) \rangle - \langle 1 | \varphi_ {2} \rangle \langle \varphi_ {1} | 1 \rangle | = 0;
\]

\begin{enumerate}
\def\labelenumi{(\roman{enumi})}
\setcounter{enumi}{2}
\item
  The map \(f \times f: X \times X \to X \times X\) is \((\mu \otimes \mu)\) -ergodic (cf. exercise 16 of I, p. 190);
\item
  The eigenspace associated with the eigenvalue 1 has dimension 1, and for every \(\lambda \in \mathbf{C} - \{1\}\) , the eigenspace of \(u_{f}\) relative to \(\lambda\) is reduced to 0.
\end{enumerate}

(To prove that (iii) implies (i), use Exercise 64 of II, p. 300.)

\begin{enumerate}
\def\labelenumi{\alph{enumi})}
\setcounter{enumi}{1}
\item
  Consider the example from part (e) of Exercise 16 in I, p. 190. The map f is weakly mixing relative to \(\mu\) .
\item
  Consider the example from question f) of exercise 16 in I, p. 190. The map f is not weakly mixing relative to \(\mu\) .
\end{enumerate}

\begin{enumerate}
\def\labelenumi{\arabic{enumi})}
\setcounter{enumi}{14}
\tightlist
\item
  Let \(G = (S, F, o, t, \iota)\) be a graph (TA, II, p. 155, def. 1). We assume that G is connected, that S is countable, and that S has at least two elements. For every \(x \in S\) , we denote by \(v(x)\) the number of edges of G of which x is an endpoint; we have \(v(x) \geqslant 1\) for all \(x \in S\) , and we assume that \(v(x)\) is finite for all \(x \in S\) . For all x and y in S, we denote by \(a(x, y)\) the number of arrows in G with endpoints x and y. We assume that \(a(x, y)\) is finite for all \((x, y) \in S^{2}\) .
\end{enumerate}

We define a measure \(\mu\) on S by \(\mu(\{x\}) = v(x)\) for all \(x \in S\) . The measure \(\mu\) is positive. We denote by \(\mathcal{F}(\mathrm{G})\) the vector space of functions on S with complex values. We denote by \(\widetilde{\mathrm{M}}_{\mathrm{G}}\) the endomorphism of \(\mathcal{F}(\mathrm{G})\) such that

\[
\widetilde {\mathrm{M}} _ {\mathrm{G}} (\varphi) (x) = \frac {1}{v (x)} \sum_ {y} a (x, y) \varphi (y),
\]

for all \(\varphi\in\mathcal{F}(\mathrm{G})\) and every \(x\in\mathrm{S}\) (“Markov operator”).

\begin{enumerate}
\def\labelenumi{\alph{enumi})}
\item
  The endomorphism \(\widetilde{\mathrm{M}}_{\mathrm{G}}\) induces, by passage to subspaces, an endomorphism \(\mathrm{M_G}\) of \(\mathrm{L}_{\mathbf{C}}^{2}(\mathrm{S},\mu)\) .
\item
  The endomorphism \(M_{G}\) is self-adjoint. It satisfies
\end{enumerate}

\[
\langle \varphi | 1 - \mathrm{M} _ {\mathrm{G}} (\varphi) \rangle = \frac {1}{2} \sum_ {(x, y) \in \mathrm{S} ^ {2}} a (x, y) | \varphi (x) - \varphi (y) | ^ {2}
\]

\[
\langle \varphi | 1 + \mathrm{M} _ {\mathrm{G}} (\varphi) \rangle = \frac {1}{2} \sum_ {(x, y) \in \mathrm{S} ^ {2}} a (x, y) | \varphi (x) + \varphi (y) | ^ {2}
\]

for all \(\varphi\in\mathrm{L}^{2}(\mathrm{S},\mu)\). Moreover, we have \(-1\leqslant\mathrm{M}_{\mathrm{G}}\leqslant1\) .

\begin{enumerate}
\def\labelenumi{\alph{enumi})}
\setcounter{enumi}{2}
\tightlist
\item
  Let \(d \geqslant 1\) be an integer and suppose that G is the Cayley graph of \(Z^{d}\) relative to the set \(\{e_{i}\}\) , where \((e_{i})\) is the canonical basis of \(Z^{d}\) (TA, II, p. 221, exercise 7). The spectral radius \(\varrho(\mathrm{M}_{\mathrm{G}})\) of \(M_{G}\) is equal to 1.
\end{enumerate}

\(d\) ) Let \(d \geqslant 3\) be an integer. Suppose that G is an infinite tree (TA, II, p. 157) \(d\) -regular, that is, an infinite tree in which every vertex has \(d\) neighbors. Let \(x_0\) be a fixed vertex of G. For every \(\lambda \in \mathbf{C}\) there exists a unique function \(\varphi_{\lambda} \colon S \to \mathbf{C}\) such that \(\varphi_{\lambda}(x_0) = 1\) and \(\widetilde{\mathrm{M}}_{\mathrm{G}}(\varphi_{\lambda}) = \lambda \varphi_{\lambda}\) . Moreover, for every \(x \in S\) the value of \(\varphi_{\lambda}\) at \(x\) depends only on the distance from \(x\) to \(x_0\) (cf. TA, II, p. 222, exercise 10).

\begin{enumerate}
\def\labelenumi{\alph{enumi})}
\setcounter{enumi}{4}
\tightlist
\item
  The spectrum of \(M_{G}\) is the interval \([-2\sqrt{d-1}/d, 2\sqrt{d-1}/d]\) ("Kesten's theorem"). In particular, we have \(\varrho(\mathrm{M}_{\mathrm{G}})=2\sqrt{d-1}/d\) .
\end{enumerate}

\begin{enumerate}
\def\labelenumi{\arabic{enumi})}
\setcounter{enumi}{15}
\tightlist
\item
  We keep the notation from the previous exercise, and we assume that S and F are finite. We denote by \(h_{G}\) the Cheeger constant of the graph G, defined by
\end{enumerate}

\[
h _ {\mathrm{G}} = \inf _ {\varnothing \neq \mathrm{V} \subset \mathrm{S}} \frac {\operatorname{Card} (\mathrm{E} (\mathrm{V} , \mathrm{S} - \mathrm{V}))}{\inf (\operatorname{Card} (\mathrm{V}) , \operatorname{Card} (\mathrm{S} - \mathrm{V}))},
\]

where E(V, W) denotes the set of arrows of G having one endpoint in V and one in W.

\begin{enumerate}
\def\labelenumi{\alph{enumi})}
\tightlist
\item
  One has \(\varrho(\mathrm{M}_{\mathrm{G}})=1\) and \(\lambda=1\) is an eigenvalue of \(\mathrm{M}_{\mathrm{G}}\) of multiplicity 1. We have \(-1\in\mathrm{Sp}(\mathrm{M}_{\mathrm{G}})\) if and only if G is bipartite (cf. TA, II, p. 219, exercise 2); in this case, the spectrum of \(\mathrm{M}_{\mathrm{G}}\) is invariant under \(\lambda\mapsto-\lambda\) and the spectral multiplicity of \(\lambda\) is equal to that of \(-\lambda\) for all \(\lambda\in\mathbf{C}\) .
\end{enumerate}

We denote by \(\mathrm{L}_{0}^{2}(\mathrm{S},\nu)\) the orthogonal complement in \(\mathrm{L}_{\mathbf{C}}^{2}(\mathrm{S},\nu)\) of the union of the eigenspaces of \(M_{G}\) corresponding to the eigenvalues 1 and -1. The endomorphism \(M_{G}\) induced by passage to subspaces is an endomorphism of \(\mathrm{L}_{0}^{2}(\mathrm{S},\nu)\), denoted \(M_{G,0}\). We have \(\varrho(\mathrm{M}_{\mathrm{G},0}) < 1\). We denote by \(\lambda_{1}(\mathrm{G})\) the smallest nonzero eigenvalue of \(1 - M_{G}\).

\begin{enumerate}
\def\labelenumi{\alph{enumi})}
\setcounter{enumi}{1}
\tightlist
\item
  Let \(v_{+} = \sup v(x)\) and \(v_{-} = \inf v(x)\) . We have
\end{enumerate}

\[
\lambda_ {1} (\mathrm{G}) \leqslant \frac {2 v _ {+}}{v _ {-} ^ {2}} h _ {\mathrm{G}}.
\]

\begin{enumerate}
\def\labelenumi{\alph{enumi})}
\setcounter{enumi}{2}
\tightlist
\item
  Let \(\varphi\colon S\to\mathbf{R}\) be a nonconstant function on \(S\) with real values; let \(m=\inf\varphi(x)\) and \(M=\sup\varphi(x)\) . Let \(t_{0}\in\mathbf{R}\) be such that, for every \(t\in\mathbf{R}\) , the set \(W_{t}=\overline{\varphi}^{-1}(]-\infty,t_{0}[)\) satisfies \(\mathrm{Card}(W_{t})\leqslant\mathrm{Card}(S)/2\) if and only if \(t<t_{0}\) . Let \(\nu\) be a diffuse measure (INT, V, p. 67, def. 5) supported on \([m,M]\) . Set
\end{enumerate}

\[
\mathrm{A} = \frac {1}{2} \sum_ {(x, y) \in \mathrm{S} ^ {2}} a (x, y) \nu ([ \varphi (x), \varphi (y) ])
\]

\[
\mathrm{B} = \sum_ {x \in \mathrm{S}} \nu ([ t _ {0}, \varphi (x) ]),
\]

where, if \(a < b\) are real numbers, we denote \(\nu([b, a]) = \nu([a, b])\) . There exists \(t \in \mathbf{R}\) such that

\[
\frac {\operatorname{Card} \left(\mathrm{E} \left(\mathrm{W} _ {t} , \mathrm{S} - \mathrm{W} _ {t}\right) \right.}{\inf \left(\operatorname{Card} \left(\mathrm{W} _ {t}\right) , \operatorname{Card} \left(\mathrm{S} - \mathrm{W} _ {t}\right)\right)} \leqslant \frac {\mathrm{A}}{\mathrm{B}}.
\]

\begin{enumerate}
\def\labelenumi{\alph{enumi})}
\setcounter{enumi}{3}
\tightlist
\item
  One has
\end{enumerate}

\[
h _ {\mathrm{G}} \leqslant v _ {+} \sqrt {2 \lambda_ {1} (\mathrm{G})}.
\]

(Apply the previous question with \(\varphi\) an eigenfunction of \(\mathrm{M_G}\) for the eigenvalue \(\lambda_1(\mathrm{G})\) and with a suitably chosen measure \(\nu\).)

\begin{enumerate}
\def\labelenumi{\alph{enumi})}
\setcounter{enumi}{4}
\tightlist
\item
  One says that a family \((\mathrm{G}_{j})_{j\in\mathrm{J}}\) of finite graphs \(\mathrm{G}_{j}=(\mathrm{S}_{j},\mathrm{F}_{j},o_{j},t_{j},i_{j})\) is a family of expander graphs if the following conditions hold:
\end{enumerate}

\begin{enumerate}
\def\labelenumi{(\roman{enumi})}
\item
  One has \(\mathrm{Card}(\mathrm{S}_{j}) \to +\infty\) according to the filter of complements of finite subsets of J;
\item
  There exists \(C \geqslant 0\) such that for every j and every \(x \in S_{j}\) , the valence of \(G_{j}\) at x is \(\leqslant C\) ;
\item
  There exists \(\delta > 0\) such that \(h_{\mathrm{G}_j} \geqslant \delta\) for all \(j \in \mathrm{J}\) .
\end{enumerate}

Show that \((\mathrm{G}_{j})\) is a family of expander graphs if and only if (i), (ii) are satisfied as well as the condition

(iii') There exists \(\delta > 0\) such that \(\lambda_1(G_j) \geqslant \delta\) for all \(j \in J\) .

\begin{enumerate}
\def\labelenumi{\arabic{enumi})}
\setcounter{enumi}{16}
\tightlist
\item
  \begin{enumerate}
  \def\labelenumii{\alph{enumii})}
  \tightlist
  \item
    Let \(n \in \mathbf{N}\). Let E be a Hilbert space of dimension \(n\), let \(u \in \mathcal{L}(\mathrm{E})\), and let \(x \in \mathrm{E}\) be of norm 1 such that \(u(x) = 0\). For every vector \(y \in \mathrm{E}\) of norm 1, we denote by \(\Delta_u(y)\) the discriminant of the quadratic form defined on \(y^\circ\) by \(z \mapsto \langle z | u(z) \rangle\) .
  \end{enumerate}
\end{enumerate}

For every \(y \in E\) of norm 1, we have \(|\langle x | y \rangle|^{2} \Delta_{u}(x) = \Delta_{u}(y)\) . (Note that \(\Delta_{u}(y) = \langle a | \bigwedge u^{n-1}(a) \rangle / \langle a | a \rangle\) for every nonzero element a of the space \(\bigwedge^{n-1} y^{\circ} \subset \bigwedge^{n-1} E\) , the space \(\bigwedge^{n-1} E\) being endowed with its canonical Hilbertian structure, cf. EVT, V, p.33; then verify the formula separately when x is orthogonal to y and when x = y.)

\begin{enumerate}
\def\labelenumi{\alph{enumi})}
\setcounter{enumi}{1}
\tightlist
\item
  Let A be a matrix of type \((n,n)\) with complex coefficients such that \({}^{t}\overline{A}=A\) . For every integer i such that \(1\leqslant i\leqslant n\) , let \(A_{i}\) be the matrix of type \((n-1,n-1)\) obtained by deleting the \(i\)-th row and the \(i\)-th column of A. Let \((e_{1},\ldots,e_{n})\) be an orthonormal basis of \(C^{n}\) , endowed with its canonical Hilbertian structure, consisting of eigenvectors for A, and let \(\lambda_{i}\) be the eigenvalue of u relative to \(e_{i}\) . Let us write \(e_{i}=(e_{i,j})_{1\leqslant j\leqslant n}\) .
\end{enumerate}

For all \((i,j)\) , we have

\[
\left| e _ {i, j} \right| ^ {2} \prod_ {k \neq i} \left(\lambda_ {i} - \lambda_ {k}\right) = \det \left(\lambda_ {i} - A _ {j}\right).
\]

(Reduce to the case where \(\lambda_{i}=0\) , then apply the preceding question to the space \(C^{n}\) , to the endomorphism \(u\colon y\mapsto Ay\) , and to the vector \(x=e_{i}\) .)

\begin{enumerate}
\def\labelenumi{\arabic{enumi})}
\setcounter{enumi}{17}
\tightlist
\item
  Let X be a compact metrizable space and let \(\mu\) be a positive measure on X. Denote by \(\mathrm{E} = \mathrm{L}_{\mathbf{C}}^{2}(\mathrm{X}, \mu)\). Let \((\mathcal{X}_{n})_{n \in \mathbf{N}}\) be a sequence of finite partitions of X into measurable sets such that \(X_{n+1}\) is finer than \(X_{n}\) for all \(n \in N\), and such that the largest of the diameters of a part of \(X_{n}\) tends to 0 as \(n \to +\infty\) .
\end{enumerate}

Let \(E_{n} \subset E\) be the space of \(f \in \mathrm{L}^{2}(\mathrm{X}, \mu)\) such that f is \(\mu\)-almost everywhere constant on Y for every part Y of \(X_{n}\). We set \(E_{-1} = \{0\}\). For every \(n \geqslant -1\), let \(p_{n}\) be the orthogonal projector onto E with image \(E_{n}\) .

\begin{enumerate}
\def\labelenumi{\alph{enumi})}
\item
  The sequence \((p_{n})_{n\in\mathbf{N}}\) converges to the identity in the space \(\mathcal{L}_{s}(\mathrm{E})\) of endomorphisms of E endowed with the topology of simple convergence. (First consider the limit of \(p_{n}(f)\) as \(f\in\mathcal{C}(\mathrm{X})\) .)
\item
  For every \(n \in N\) , the endomorphism \(q_{n} = p_{n} - p_{n-1}\) is the orthogonal projector onto \(F_{n} = E_{n} \cap E_{n-1}^{\circ}\) , and the spaces \(F_{n}\) are pairwise orthogonal.
\item
  One has \(\sum_{n} q_{n} = 1_{\mathrm{E}}\) in \(\mathcal{L}_s(\mathrm{E})\) , and the space E is the Hilbertian sum of the spaces \(\mathrm{F}_n\) .
\end{enumerate}

\begin{enumerate}
\def\labelenumi{\arabic{enumi})}
\setcounter{enumi}{18}
\tightlist
\item
  Let \(\mathrm{X} \subset \mathbf{R}\) be a compact subset of \(\mathbf{R}\) and let \(\mu\) be a positive measure on \(\mathrm{X}\) and let \(u\) be the endomorphism of multiplication by the identity function of \(\mathrm{X}\) .
\end{enumerate}

We denote \(\mathrm{E}=\mathrm{L}^{2}(\mathrm{X},\mu)\) and we denote by \(\mathcal{L}_{s}(\mathrm{E})\) the space of endomorphisms of E endowed with the topology of simple convergence.

Let \(\varepsilon > 0\). Let \(r > 0\) such that \(\mathrm{X} \subset ] - r, r[\) and let \(\ell \geqslant 1\) such that \(r2^{-\ell} < \varepsilon\) . For every integer \(n \in \mathbf{N}\) and every integer \(j\) such that \(|j| \leqslant 2^{n + \ell}\) , let

\[
\mathrm{X} _ {n, j} = \mathrm{X} \cap \Big ] \frac {r j}{2 ^ {n + \ell}}, \frac {r (j + 1)}{2 ^ {n + \ell}} \Big ]
\]

\begin{enumerate}
\def\labelenumi{\alph{enumi})}
\item
  For every integer \(n \in N\), the family \(\mathcal{X}_{n} = (\mathrm{X}_{n,j})_{j}\) is a finite partition of X; the sets \(X_{n,j}\) are measurable and of diameter \(\leqslant r2^{-n-\ell}\) . For every \(n \in N\), the partition \(X_{n+1}\) is finer than the partition \(X_{n}\). We can therefore apply to \((\mathcal{X}_{n})\) the results of the previous exercise, whose notation we reuse, in particular the spaces \(E_{n}\) and \(F_{n}\) and the orthogonal projectors \(p_{n}\) and \(q_{n}\) . The spaces \(E_{n}\) are finite-dimensional.
\item
  The series
\end{enumerate}

\[
v = \sum_ {n \in \mathbf {N}} q _ {n} u q _ {n}
\]

converges in the space \(\mathcal{L}_{s}(\mathrm{E})\); there exists an orthonormal basis B of E such that \(v \in \mathcal{D}_{\mathrm{B}}(\mathrm{E})\) .

\begin{enumerate}
\def\labelenumi{\alph{enumi})}
\setcounter{enumi}{2}
\tightlist
\item
  The endomorphism \(w = v - u\) is compact. (Consider the endomorphisms \(u_n\) of multiplication by the functions \(g_n \in \mathrm{E}_n\) such that, for every integer \(j\) such that \(|j| \leqslant 2^{n + \ell}\) , and for all \(x \in \mathrm{X}_{n,j}\) , we have \(g_n(x) = rj2^{-n - \ell}\) ; then verify that \(\| u - u_n \| \leqslant r2^{-n - \ell}\) , and such that \(u_n p_n = p_n u_n\) .)
\end{enumerate}

\begin{enumerate}
\def\labelenumi{\arabic{enumi})}
\setcounter{enumi}{19}
\tightlist
\item
  Let E be a complex Hilbert space of countable type and \(u \in \mathcal{L}(\mathrm{E})\) be a Hermitian endomorphism. Let \(\varepsilon\) be a strictly positive real number. There exists an orthonormal basis B of E, a Hermitian diagonal endomorphism \(v \in \mathcal{D}_{\mathrm{B}}(\mathrm{E})\) and a compact endomorphism \(w \in \mathcal{L}^{\mathrm{c}}(\mathrm{E})\) such that \(u = v + w\) and \(\|w\| < \varepsilon\) (First consider the case where u admits a cyclic vector, and then apply the preceding exercise.)
\end{enumerate}

21) Let E be a complex Hilbert space of countable type and \(u \in \mathcal{L}(\mathrm{E})\) a normal endomorphism. Let \(\varepsilon\) be a strictly positive real number. There exists an orthonormal basis B of E, a diagonal endomorphism \(v \in \mathcal{D}_{\mathrm{B}}(\mathrm{E})\) and a compact endomorphism \(w \in \mathcal{L}^{\mathrm{c}}(\mathrm{E})\) such that \(u = v + w\) and \(\|w\| < \varepsilon\) .

\begin{enumerate}
\def\labelenumi{\arabic{enumi})}
\setcounter{enumi}{21}
\tightlist
\item
  Let G be a locally compact commutative topological group and \(\mu\) a Haar measure on G.
\end{enumerate}

Let \(u \in \mathcal{L}(\mathrm{L}^2(\mathrm{G},\mu))\) be such that \(u\) commutes with the multiplication operator \(m_{\chi}\) for every character \(\chi \in \widehat{\mathrm{G}}\). Prove that there exists \(g \in \mathrm{L}^{\infty}(\mathrm{G},\mu)\) such that \(u = m_g\) .

\begin{enumerate}
\def\labelenumi{\arabic{enumi})}
\setcounter{enumi}{22}
\tightlist
\item
  Let \(u\) be an endomorphism of E and let A be the closed unital subalgebra of \(\mathcal{L}(\mathrm{E})\) generated by \(u\) .
\end{enumerate}

If v is an endomorphism of E such that \(u \circ v - v \circ u\) belongs to A, then \(u \circ v - v \circ u\) is quasi-nilpotent. In particular, we have \(u \circ v - v \circ u \neq 1_{E}\) if E is nonzero. (Use exercise 10 of I, p. 167; note that the result is valid for any Banach space E.)

\begin{enumerate}
\def\labelenumi{\arabic{enumi})}
\setcounter{enumi}{23}
\tightlist
\item
  Let \(\mathrm{A} \subset \mathcal{L}(\mathrm{E})\) be a unital self-adjoint subalgebra of \(\mathcal{L}(\mathrm{E})\) . We denote by \(\mathrm{A}''\) the bicommutant of A. We denote by \(\mathcal{L}(\mathrm{E})_s\) the space \(\mathcal{L}(\mathrm{E})\) endowed with the topology of simple convergence.
\end{enumerate}

\begin{enumerate}
\def\labelenumi{\alph{enumi})}
\item
  The space \(A''\) is a unital self-adjoint subalgebra of \(\mathcal{L}(\mathrm{E})\) which is closed in \(\mathcal{L}(\mathrm{E})_s\) .
\item
  Let \(v \in A''\) and \(x \in E\). The vector \(v(x)\) belongs to the closure in E of the set of \(u(x)\) for \(u \in A\). (Let F be this closure; observe that the orthogonal projector of E with image F belongs to \(A'\).)
\item
  The algebra \(A''\) coincides with the closure of A in \(\mathcal{L}(\mathrm{E})_s\) . (For every finite family \((x_i)_{i \in \mathrm{I}}\) of elements of E, apply the preceding question to the space \(\mathrm{E}^{\mathrm{I}}\) , to the image \(\widetilde{\mathrm{A}}\) of A by the unital involutive morphism \(f: \mathcal{L}(\mathrm{E}) \to \mathcal{L}(\mathrm{E}^{\mathrm{I}})\) such that \(f(u)(y_i) = (u(y_i))_{i \in \mathrm{I}}\) and to the element \((x_i)\) of \(\mathrm{E}^{\mathrm{I}}\) .)
\end{enumerate}

The equality of \(A''\) and the closure of A in \(\mathcal{L}(\mathrm{E})_s\) is von Neumann's bicommutant theorem.

\begin{enumerate}
\def\labelenumi{\arabic{enumi})}
\setcounter{enumi}{24}
\tightlist
\item
  Let U be an open subset of R. A universally measurable function \(f: U \to R\) is called a Loewner function if, for every complex Hilbert space E and all Hermitian endomorphisms u and v of E such that \(\operatorname{Sp}(u) \subset \operatorname{U}\) and \(\operatorname{Sp}(v) \subset \operatorname{U}\) , the condition \(u \leqslant v\) in \(\mathcal{L}(\operatorname{E})\) implies \(f(u) \leqslant f(v)\) . We denote by \(\operatorname{Loe}(\operatorname{U})\) the set of Loewner functions on U.
\end{enumerate}

\begin{enumerate}
\def\labelenumi{\alph{enumi})}
\item
  Let \(f \in \text{Loe}(U)\). The function f is increasing on U. The constant function 1 and the identity function of U belong to \(\text{Loe}(U)\) .
\item
  Let \(f \in \text{Loe}(\mathbf{R}_+^*)\). Let E be a complex Hilbert space and let u be a positive endomorphism of E. For every orthogonal projector p of E and for every \(\varepsilon > 0\) , we have \(f(u)_p \leqslant f((1 + \varepsilon)u_p)\) , where we denote by \(v_p\) the endomorphism of \(\text{Im}(p)\) defined by \(x \mapsto p(v(x))\) for all \(v \in \mathcal{L}(\text{E})\) . (Identify E with the Hilbert direct sum \(\text{Im}(p) \oplus \text{Im}(1_{\text{E}} - p)\) and represent u as a matrix \(\begin{pmatrix} a & b \\ c & d \end{pmatrix}\). Note that \(u \leqslant \begin{pmatrix} (1 + \varepsilon)a & 0 \\ 0 & (1 + \varepsilon^{-1})d \end{pmatrix}\) for all \(\varepsilon > 0\) .)
\item
  Let \(f \in \text{Loe}(\mathbf{R}_{+}^{*})\). The function f is concave on \(R_{+}^{*}\) (Let v be an endomorphism of \(C^{2}\) diagonal in the canonical basis and let \(g \in \text{SO}(\mathbf{R}^{2})\) viewed as an endomorphism of \(C^{2}\). Apply the previous question to the orthogonal projector with image \(Ce_{1}\) and with \(u = g^{*}vg\) .)
\item
  The set \(\operatorname{Loe}(\mathbf{R})\) is the set of functions of the form \(f(x)=ax+b\) where \((a,b)\in\mathbf{R}_{+}\times\mathbf{R}\) . (Apply the previous question to f as well as to the function \(x\mapsto-f(-x)\) .)
\item
  Let \(t \in \mathbf{R}_{+}^{*}\) . The function \(f: x \mapsto -(t + x)^{-1}\) belongs to \(\operatorname{Loe}(\mathbf{R}_{+}^{*})\).
\end{enumerate}

\(f)\) For any bounded positive measure \(\nu\) on \(\mathbf{R}_{+}^{*}\) such that the function \(x\mapsto x^{-1}\) is \(\nu\) -integrable, the function \(f\colon\mathbf{R}_{+}^{*}\to\mathbf{R}\) defined by

\[
f (x) = \int_ {\mathbf {R} _ {+} ^ {*}} \frac {x}{x + t} d \nu (t)
\]

belongs to \(\text{Loe}(\mathbf{R}_{+}^{*})\); it is differentiable on \(R_{+}^{*}\) and moreover satisfies the following conditions:

\begin{enumerate}
\def\labelenumi{(\roman{enumi})}
\item
  \(\lim_{t\to0}f(t)=0;\)
\item
  \(\lim_{t\to0}f'(t)\) exists in R;
\item
  f has a finite limit at \(+\infty\) .
\end{enumerate}

\begin{enumerate}
\def\labelenumi{\arabic{enumi})}
\setcounter{enumi}{25}
\tightlist
\item
  Let E be a complex Hilbert space and let a be a positive endomorphism of E. We denote by \(\|x\|_{a} = \langle x | a(x) \rangle\) ; it is a seminorm on E. We likewise denote by \(\|u\|_{a}\) the seminorm of an endomorphism u of the space E equipped with this seminorm.
\end{enumerate}

We call a function \(f \colon \mathbf{R}_+^* \to \mathbf{R}_+^*\) an interpolation function if, for every complex Hilbert space E, every positive endomorphism \(a\) of E and every \(u \in \mathcal{L}(\mathrm{E})\) , the conditions \(\| u \| \leqslant 1\) and \(\| u \|_a \leqslant 1\) imply that \(\| f(u) \|_a \leqslant 1\) .

\begin{enumerate}
\def\labelenumi{\alph{enumi})}
\item
  A function \(f \colon \mathbf{R}_+^* \to \mathbf{R}_+^*\) is an interpolation function if and only if, for every complex Hilbert space E, every positive endomorphism \(a\) of E and every \(u \in \mathcal{L}(\mathrm{E})\) , the conditions \(u^*u \leqslant 1_{\mathrm{E}}\) and \(u^*au \leqslant a\) imply \(u^*f(a)u \leqslant f(a)\) .
\item
  Any Loewner function on \(R_{+}^{*}\) is an interpolation function. (Given E and \(u \in \mathcal{L}(E)\) of norm \(\leqslant 1\), set \(w = \sqrt{1_{E} - u^{*}u}\) and \(w' = \sqrt{1_{E}-uu^{*}}\); consider \(F = E \oplus E\) and note that the endomorphism v of F defined by \(v(x,y) = (u(x)+w'(y), -w(x)+u(y))\) is an isometry; then consider \(a' \colon (x,y) \mapsto (a(x),0)\) and apply part b) of the previous exercise to the orthogonal projector with image \(E \oplus \{0\}\) .)
\item
  Every interpolation function is a Loewner function on \(\mathbf{R}_{+}^{*}\) . (Let \(0 \leqslant u \leqslant v\) be endomorphisms of a complex Hilbert space E with \(v\) invertible; set \(\mathrm{F} = \mathrm{E} \oplus \mathrm{E}\) and consider the endomorphism \(a = u \oplus v\) of F as well as the endomorphism \(w\) represented by the matrix \(\left( \begin{array}{cc}0 & 1_{\mathrm{E}}\\ v^{-1/2}u^{1/2} & 0 \end{array} \right)\) ; show that \(w^{*}w \leqslant 1_{\mathrm{F}}\) and \(w^{*}aw \leqslant a\) .)
\item
  Let f be an interpolation function. Let E be a complex Hilbert space; we denote \(\mathcal{R}(u)=\frac{1}{2}(u+u^{*})\) for all \(u\in\mathcal{L}(\mathrm{E})\) . For u and v in \(\mathcal{L}(\mathrm{E})\) , prove that \(e^{-tv^{*}}\circ u\circ e^{-tv}\leqslant u\) for all \(t\in R_{+}\) if and only if \(\mathcal{R}(u\circ v)\geqslant0\) . (For a given \(x\in E\), compute the derivative of the function \(t\mapsto\langle e^{-tv}x|(ue^{-tv})x\rangle\) .)
\end{enumerate}

Hence deduce that for all endomorphisms \(u\) and \(v\) of E such that the conditions

\[
u \geqslant 0, \quad \mathcal {R} (v) \geqslant 0, \quad \mathcal {R} (u v) \geqslant 0
\]

are satisfied, we have \(\mathcal{R}(f(u)v) \geqslant 0\) .

27) We propose to prove the converse of the assertion in question \(b\) of Exercise 25. For this, one will use the results and notations of Exercise 23 of V, p. 503 below. In particular, G denotes the group \(\mathbf{R}_{+}^{*}\) equipped with the Haar measure \(\mu = x^{-1}dx\) in duality with the group \(i\mathbf{R}\) by the mapping \((x, it) \mapsto x^{it}\) . We denote by X the identity map of G.

We denote by co, si, and ta the functions \(it \mapsto \cos(i\pi t)\), \(it \mapsto \sin(i\pi t)\), and \(it \mapsto \tan(i\pi t)\) on iR, respectively.

Let \(f \in \text{Loe}(\mathbf{G})\) satisfy conditions (i), (ii), and (iii) of question \(f\) of Exercise 25.

\begin{enumerate}
\def\labelenumi{\alph{enumi})}
\item
  Let \(\alpha\) be the function \(x \mapsto x / (1 + x)\) on \(G\). We have \(\mathrm{I}_{\alpha} = ] - 1,0[\) and \(\widehat{\alpha}(s) = -\pi / \sin(\pi s)\) for \(s \in \mathrm{B}_f\) .
\item
  The function ta is bounded on \(i\mathbf{R}\). We set \(v = F \circ m_{ta} \circ \overline{F}\) ; it is an endomorphism of \(L^{2}(G)\) such that \(\|v\| \leqslant 1\) and \(\mathcal{R}(v) = 0\) .
\item
  Let \(\varepsilon > 0\) be a real number. We denote by \(u_{\varepsilon}\) the multiplication endomorphism by \(h_{\varepsilon} = \mathrm{X}/(1 + \varepsilon\mathrm{X})\) on \(\mathrm{L}^2(\mathrm{G})\). We have \(u_{\varepsilon} \geqslant 0\) and for every \(\varphi \in \mathrm{L}^2(\mathrm{G})\) and every bounded continuous function \(f\) on \(\mathrm{G}\) , we have
\end{enumerate}

\[
\langle \varphi | \mathcal {R} (f (u _ {\varepsilon}) v) \varphi \rangle = \frac {1}{2} \langle f \circ h _ {\varepsilon} | \mathrm{Q} (\varphi) \rangle
\]

\(\mathrm{Q}(\varphi) = \overline{\varphi} v(\varphi) + \overline{v(\varphi)}\varphi .\)

\begin{enumerate}
\def\labelenumi{\alph{enumi})}
\setcounter{enumi}{3}
\item
  Let F be the vector subspace of \(\mathrm{L}^{2}(i\mathbf{R})\) generated by the functions \(t \mapsto e^{\alpha(it - s_{0})^{2}}\) where \(\alpha \in R_{+}^{*}\) and \(s_{0} \in C\); it is dense in \(\mathrm{L}^{2}(i\mathbf{R})\) .
\item
  Let \(\psi \in \mathrm{F}\) . There exists \(\varphi \in \mathrm{L}^2 (\mathbf{G})\) such that \(\widehat{\varphi} = \mathrm{co}\cdot \psi\). Let \(g = \mathrm{Q}(\varphi)\). We have \(g\in \mathrm{L}^{1}(\mathrm{G})\) , and there exists a holomorphic function on \(\mathbf{C}\) whose restriction to \(i\mathbf{R}\) coincides with \(\widehat{g}\) . Moreover
\end{enumerate}

\[
\widehat {g} = \frac {1}{2 i \pi} (\psi^ {*} * \psi) \mathrm{si}.
\]

\(f)\) For every \(\varepsilon > 0\) , we have \(\mathcal{R}(u_{\varepsilon} \circ v) \geqslant 0\) . (Compute \(\langle \varphi | \mathcal{R}(u_{\varepsilon} u) \varphi \rangle\) when \(\varphi\) is a function as in the preceding question, and use convolution on G as in question \(d\) of Exercise 23 of V, p. 503.)

\(g)\) Let \(f \in \mathrm{Loe}(\mathbf{R}_{+}^{*})\). Let \(g\) be a function associated with \(\psi \in F\) as in question \(e)\). We then have

\[
\int_ {\mathrm{G}} f g d \mu \geqslant 0
\]

and

\[
\int_ {\mathbf {R}} \widehat {f} (- \sigma - i t) \sin (\pi (\sigma + i t)) (\psi^ {*} * \psi) (\sigma + i t) d t \geqslant 0
\]

for all \(\sigma\in]0,1[\) .

\begin{enumerate}
\def\labelenumi{\alph{enumi})}
\setcounter{enumi}{7}
\tightlist
\item
  Let \(\sigma \in ]0,1[\) . The function \(q\) defined by \(q(s) = \widehat{f}(-s)\sin(\pi s)\) for \(s\) such that \(\mathcal{R}(s) = \sigma\) satisfies
\end{enumerate}

\[
\sum_ {i = 0} ^ {n} \sum_ {j = 0} ^ {n} \bar {t} _ {i} t _ {j} q \left(\frac {1}{2} (\bar {z} _ {i} + z _ {j})\right) \geqslant 0
\]

for all \(n \in N\) and all families \((t_{i})_{0 \leqslant i \leqslant n}\) and \((z_{i})_{0 \leqslant i \leqslant n}\) of complex numbers such that \(\mathcal{R}(z_{i}) = \sigma\) for every i.

\begin{enumerate}
\def\labelenumi{\roman{enumi})}
\tightlist
\item
  There exists a positive measure \(\nu_0\) on G such that \(\widehat{\nu}_0(it) = \widehat{f}(-it)\sin(i\pi t)/\pi\) for all \(t\in\mathbf{R}\) ; we have \(\mathrm{X}^{\sigma}\in\mathrm{L}^{1}(\mathrm{G},\nu_{0})\) if \(0<\sigma<1\) .
\end{enumerate}

\begin{enumerate}
\def\labelenumi{\alph{enumi})}
\setcounter{enumi}{9}
\tightlist
\item
  Let \(\nu = \check{\nu}_0\). The measure \(\nu\) is bounded. (Verify that \(]-1,0[ \subset I_{\nu}\) and then that \(\widehat{\nu}(-t)\) is bounded when \(t \to 0\) to show that \(\nu([1,+\infty[)\) is finite.) For every \(x\) , we have
\end{enumerate}

\[
f (x) = \int_ {\mathrm{G}} \frac {x}{x + t} d \nu (t).
\]

\begin{enumerate}
\def\labelenumi{\arabic{enumi})}
\setcounter{enumi}{27}
\tightlist
\item
  \begin{enumerate}
  \def\labelenumii{\alph{enumii})}
  \tightlist
  \item
    Let \(D \subset C\) be the open unit disk. Let f be a germ of a holomorphic function in a neighborhood of \(\overline{D}\) . For every \(z \in D\) , we have
  \end{enumerate}
\end{enumerate}

\[
f (z) = i \mathcal {I} (f (0)) + \frac {1}{2 \pi} \int_ {0} ^ {2 \pi} \mathcal {R} (f (e ^ {i \theta})) k (e ^ {i \theta}, z) d \theta
\]

where \(k(z, w) = (z + w) / (z - w)\) for \((z, w) \in \mathrm{D} \times \mathrm{D}\) . (Note that

\[
k (e ^ {i \theta}, z) = 1 + 2 \sum_ {n = 1} ^ {+ \infty} z ^ {n} e ^ {- i n \theta}
\]

and use the Taylor expansion of f around 0.)

\begin{enumerate}
\def\labelenumi{\alph{enumi})}
\setcounter{enumi}{1}
\tightlist
\item
  Let \(f\) be a holomorphic function on D such that \(\mathcal{R}(f(z)) > 0\) for all \(z \in \mathrm{D}\) . The measures \(\mu_r = \frac{1}{2\pi} \mathcal{R}(f(re^{i\theta})) d\theta\) on \([0, 2\pi]\) converge vaguely as \(r \to 1\) to a bounded positive measure \(\nu\) , and we have
\end{enumerate}

\[
f (z) = i \mathcal {I} (f (0)) + \frac {1}{2 \pi} \int_ {0} ^ {2 \pi} \mathcal {R} (f (e ^ {i \theta})) k (e ^ {i \theta}, z) d \nu (\theta)
\]

for all \(z \in D\) . (Reduce to the case where \(f(0) = 1\) and observe that the measures \(\mu_{r}\) are then positive measures of total mass 1.)

\begin{enumerate}
\def\labelenumi{\alph{enumi})}
\setcounter{enumi}{2}
\tightlist
\item
  Let H be the open subset of C consisting of \(z \in C\) such that \(\mathcal{I}(z) > 0\). Let \(f: H \to C\) be a holomorphic function. The image of f is contained in H if and only if there exists a bounded positive measure \(\nu\) on R and a real number \(a \geqslant 0\) such that
\end{enumerate}

\[
f (z) = \mathcal {R} (f (i)) + a z + \int_ {\mathbf {R}} \frac {1 + x z}{x - z} d \nu (x)
\]

for all \(z \in \mathbf{H}\) .

\begin{enumerate}
\def\labelenumi{\alph{enumi})}
\setcounter{enumi}{3}
\tightlist
\item
  Let I be a nonempty open interval in \(\mathbf{R}\) and let \(f\colon \mathrm{I}\to \mathbf{R}\) . Denote by \(\mathrm{J} = \mathbf{R} - \mathrm{I}\) . The following conditions are equivalent («Loewner's theorem»):
\end{enumerate}

\begin{enumerate}
\def\labelenumi{(\roman{enumi})}
\item
  The function \(f\) belongs to \(\operatorname{Loe}(\mathbf{I})\) ;
\item
  There exists an analytic function \(\widetilde{f}\) on \((\mathbf{C} - \mathbf{R}) \cup I\) whose restriction to I coincides with \(f\) and which satisfies \(\mathcal{I}(f(z)) > 0\) if \(\mathcal{I}(z) > 0\) ;
\item
  There exists a bounded measure \(\nu\) on J and a pair \((a,b)\in\mathbf{R}_{+}\times\mathbf{R}\) such that
\end{enumerate}

\[
f (x) = a x + b + \int_ {\mathrm{J}} \frac {1 + x y}{y - x} d \nu (y)
\]

for all \(x \in I\) .

(First verify directly that (iii) implies (i) and (ii); to show that (i) implies (iii), treat separately the case where I = R and reduce the other cases to the case where I = R \(_{+}^{*}\) , for which one applies the preceding exercise; finally, to show that (ii) implies (iii), apply part c) and verify that the measure \(\nu\) on R given by it is supported on J.)

